% Task-oriented recipes.
\section{Recipes}\label{ch:recipes}

The tutorial introduced the foundational grammar of \pkg{}; this chapter gathers
a collection of practical recipes for the everyday situations you will encounter
in research papers.  Each recipe focuses on a single semantic task.  When
designing a diagram of your own, start from the tutorial example that most
closely matches your mathematical structure, and add only the specific policies
your physics requires.

\subsection{Keep a map open; close a scalar}

In \pkg{}, a perimeter side exposes indices only when you explicitly ask it to.
A diagram that states no boundary policy draws no dangling legs and carries an
empty boundary signature on those faces.  Consequently, an operator map keeps
its outer indices by stating \tnkey{boundary=open}---or \tnkey{west=open} for a
single side---while \tnkey{boundary=none} is simply the explicit spelling of the
default behavior.  The layout grid places the tensor atoms; it does not assume
that every perimeter cell owns a dangling index.  That is why a circuit pyramid
or a single four-legged junction never sprouts unwanted legs you did not ask for.

The rest of the boundary alphabet is equally intuitive: use
\tnkey{boundary=periodic} to trace a row back into itself, and use
\tnkey{west=cup} or \tnkey{east=cup} to connect adjacent layers.  Notice that
two closures may happen to have the same number of open ends and yet represent
entirely different mathematical operations!

\subsection{Stack ket, operator, and bra}

\begin{tnexample}[
  formula={\sum_{s,s'}A^s\otimes O^{ss'}\otimes\bar A^{s'}},
  caption={an expectation-value window},
  index={The rows declare their roles.  Vertical connections contract physical
    indices; six horizontal stubs keep the window a map.},
  signature={kernel-boundary|signature=open:e, open:e, open:e, open:w, open:w, open:w}]
% Formula: sum_{s,s'} A^s tensor O^{ss'} tensor conjugate(A^{s'}).
% Labels: automatic placement stands each name on its atom's free face --
% north for the ket row, south for the bra row.
\begin{tenkz}[rows={ket,op,bra}, cols=2, west=open, east=open]
  \tn{A} & \tn{A} \\
  \tn[skin=mpo]{O} & \tn[skin=mpo]{O} \\
  \tn{$\overline{A}$} & \tn{$\overline{A}$}
\end{tenkz}
\end{tnexample}

\subsection{Block a run without changing topology}

When coarse-graining, a bracket mark over a cell range (\tncmd{tnmark} with
\tnkey{form=bracket}) declares a brace or measured grouping over existing
atoms, while \tnval{skin=dots} declares an ellipsis cell whose wires pass
straight through.  Neither construct alters the underlying contractions!  A wide
single atom is appropriate only after your formula has already replaced the run
by a single blocked tensor.

\subsection{Distinguish a role from a species}

A good diagram distinguishes between structural roles and persistent species.
Use \tnkey{role=} for an atom's fixed semantic role, such as \tnval{operator} or
\tnval{marked}.  By contrast, declare \tnkey{species=} at setup scope when your
paper features a recurring family of tensors whose members should maintain a
consistent hue across multiple figures.  Always remember: labels distinguish
objects of the same kind; color should never be the sole carrier of identity.
A role may modify visual style, but never network topology.

\subsection{Stand a label on an open leg}

An open leg is a wire like any other, and it carries an accessible name.  The
boundary policy names its legs after the side and the row or column they
depart from---\tnval{open-west-1}, \tnval{open-south-col-3}---while the
\tnkey{physical=} policy names each leg after its face and home
cell---\tnval{leg-n-1-1} and \tnval{leg-s-2-3}.  Similarly, a user-authored
wire takes whatever name its \tncmd{tnwire} declares.  You can address a
fractional station along that named wire to place a label or bead with complete
precision:

\begin{Verbatim}[fontsize=\small,frame=leftline]
\begin{tenkz}[rows={wire}, cols=2, west=open, east=open]
  \tn{A} & \tn{B}
  \tnmark[form=label, label pos=90]{on open-west-1 0.5}{$\alpha$}
\end{tenkz}
\end{Verbatim}

The fractional station is measured along the path the wire follows.  A leg has
one end anchored to a glyph and one free tip; fraction zero is always the end
where the wire began.  A policy leg runs outward from the glyph, so \tnval{on
leg-n-1-1 0} lies on the north face of the site in row~1, column~1, while
\tnval{on leg-n-1-1 1} reaches its outer tip.  Declaring a row's legs once with
\tnkey{physical=} is therefore especially convenient: the very same name
serves both for line crossings and for label placement.

\subsection{Say which bond is the one under discussion}

A lattice draws its bonds automatically from its row and column policies, and
every one of them is typeset as a standard black contraction.  When a figure
needs to highlight one specific edge---the bond being cut, or the link being
entangled---you can restate that bond directly in the body.  An index wire
declared between two adjacent cells becomes that pair's bond: \pkg{} creates no
duplicate wire beneath it, and the wire's species, stroke, and route become the
bond's own:

\begin{Verbatim}[fontsize=\small,frame=leftline]
\tndeclare{species}{blocked-edge}{hue=source:red!65!black}
\begin{tenkz}[rows={wire,wire}, cols=3]
  \tn{} & \tn{} & \tn{}\\
  \tn{} & \tn{} & \tn{}
  \tnwire[species=blocked-edge]{(1,1)}{(1,2)}
\end{tenkz}
\end{Verbatim}

Because this restatement represents the same mathematical contraction, the
count of internal bonds and the external boundary signature remain unchanged;
only the visual ink reflects your styling.  You can address the pair by cell
coordinates as above, or by typed ports (\tnval{a.0} and \tnval{b.180}) on
named atoms.  (Specifying \tnkey{bonds=none} and redrawing every bond by hand is
needlessly tedious and completely unnecessary!)  A travelling string is not a
restatement: it carries an operator index, so it is drawn cleanly over the
lattice without retiring the underlying bonds.

\subsection{Send a physical index sideways}

In \pkg{}, an index type belongs to the port that carries it.  The \tnkey{ports=}
key types each perimeter face individually, so \tnval{0:physical} directs a
physical index eastward just as naturally as \tnval{0:virtual} directs a
virtual one.  A single-atom linear map types its opposing faces accordingly:

\begin{Verbatim}[fontsize=\small,frame=leftline]
\begin{tenkz}[rows={wire}, cols=1, bonds=none]
  \tn[at=(1,1), skin=box, ports={180:virtual:V, 0:physical:W}]{A}
\end{tenkz}
\end{Verbatim}

The \tnkey{physical=} key is a picture-wide policy, not the sole way to create
physical legs: it equips every cell atom with an outward physical port in the
frame's natural direction.  Both mechanisms complement each other harmoniously.
Picture policy shines on uniform rows: declare it once, and every site receives
its physical leg, with any rare exception opting out via \tnkey{physical=none}.
Conversely, on a sparse row where only a single site among many carries a
physical index, it is far cleaner to declare the port on that individual atom:

\begin{Verbatim}[fontsize=\small,frame=leftline]
\begin{tenkz}[rows={wire}, cols=5]
  \tn{} & \tn{} & \tn[ports={180:virtual, 0:virtual, 90:physical:$i$}]{} &
  \tn{} & \tn{}
\end{tenkz}
\end{Verbatim}

\subsection{Keep a physical index out of a projected sheet}

A PEPS tensor has five index directions, but only four lie in the plane of the
lattice sheet.  The four numeric ports below follow the projected planar axes,
while the picture policy supplies the upward transverse direction:

\begin{Verbatim}[fontsize=\small,frame=leftline]
\begin{tenkz}[rows={wire,wire,wire}, cols=3,
               frame=plane, bonds=none, physical=up]
  \tn[at=(2,2),
      ports={0:virtual,90:virtual,180:virtual,270:virtual}]{A}
\end{tenkz}
\end{Verbatim}

Do not attempt to add an artificial fifth numeric angle to straighten the
transverse leg!  A numeric angle always lives within the plane of the sheet,
whereas \tnkey{physical=up} follows the plane's independent transverse normal.

Remember also that physical indices belong to lattice sites defined by frame
addresses.  A mark positioned at a midpoint, along a wire, or at a crossing is a
geometric annotation, not a lattice site; hence, picture policies never sprout
physical legs from marks.  If an annotation represents an additional tensor with
its own physical index, simply declare that port explicitly via \tnkey{ports=}.

\subsection{Draw fusion trees; leave commutative diagrams to tikz-cd}

Use \tncmd{tntree} when you need to draw parenthesized fusion words and tree
contractions.  Commutative diagrams, by contrast, compose morphisms between
mathematical objects rather than contracting tensor indices; they belong to the
excellent \textsf{tikz-cd} package, outside the scope of \pkg{}.  A
\tncmd{tntree} atom is a self-contained box that embeds gracefully inside any
\textsf{tikz-cd} cell whenever the two worlds meet.

\subsection{Put a picture in a sentence}

A \pkg{} picture enters surrounding \TeX{} as a displayed mathematical term.
There is no special ``inline'' key: inline mathematical alignment belongs to the
ambient equation or text scope, not to an internal property of individual
diagram records.
